% !TeX program = xelatex
\documentclass[11pt,a4paper]{article}

% -------------------------------------------------
% Packages
% -------------------------------------------------
\usepackage{kotex}
\usepackage{fontspec}
\usepackage{amsmath,amssymb,bm,mathtools}
\usepackage[left=24mm,right=24mm,top=23mm,bottom=25mm]{geometry}
\usepackage{booktabs,array}
\usepackage{enumitem}
\usepackage{setspace}
\usepackage{xcolor}
\usepackage{colortbl}
\usepackage{tikz}
\usepackage{pifont}
\usetikzlibrary{positioning,arrows.meta,calc,fit,decorations.pathreplacing}
\usepackage[most]{tcolorbox}
\usepackage{listings}
\usepackage{hyperref}
\usepackage{fancyhdr}

\setstretch{1.14}
\setlength{\parindent}{0pt}
\setlength{\parskip}{0.55em}

% -------------------------------------------------
% Colors and boxes
% -------------------------------------------------
\definecolor{deepblue}{RGB}{34,71,120}
\definecolor{softblue}{RGB}{239,246,253}
\definecolor{softgray}{RGB}{246,247,249}
\definecolor{softgreen}{RGB}{239,249,242}
\definecolor{softyellow}{RGB}{255,249,229}
\definecolor{codegray}{RGB}{248,248,248}

\newtcolorbox{keybox}[1][]{
  colback=softblue,
  colframe=deepblue,
  boxrule=0.7pt,
  arc=2mm,
  title=#1,
  fonttitle=\bfseries
}

\newtcolorbox{thinkbox}[1][]{
  colback=softyellow,
  colframe=orange!70!black,
  boxrule=0.6pt,
  arc=2mm,
  title=#1,
  fonttitle=\bfseries
}

\newtcolorbox{remarkbox}[1][]{
  colback=softgray,
  colframe=black!45,
  boxrule=0.5pt,
  arc=2mm,
  title=#1,
  fonttitle=\bfseries
}

\newtcolorbox{mlbox}[1][]{
  colback=softgreen,
  colframe=green!45!black,
  boxrule=0.6pt,
  arc=2mm,
  title=#1,
  fonttitle=\bfseries
}

% -------------------------------------------------
% Code style
% -------------------------------------------------
\lstdefinestyle{pythonstyle}{
  language=Python,
  backgroundcolor=\color{codegray},
  basicstyle=\ttfamily\small,
  keywordstyle=\bfseries,
  commentstyle=\itshape\color{black!60},
  stringstyle=\color{black!75},
  numbers=left,
  numberstyle=\tiny\color{black!55},
  numbersep=8pt,
  frame=single,
  rulecolor=\color{black!25},
  breaklines=true,
  showstringspaces=false,
  tabsize=4,
  columns=fullflexible,
  keepspaces=true
}
\lstset{style=pythonstyle}

% -------------------------------------------------
% Math commands
% -------------------------------------------------
\newcommand{\R}{\mathbb{R}}
\newcommand{\grad}{\nabla}
\newcommand{\vect}[1]{\bm{#1}}
\newcommand{\T}{^{\mathsf T}}
\newcommand{\loss}{\mathcal{L}}
\newcommand{\sigmoid}{\sigma}
\newcommand{\pd}[2]{\frac{\partial #1}{\partial #2}}
\newcommand{\softmax}{\operatorname{softmax}}

% -------------------------------------------------
% Header / footer
% -------------------------------------------------
\pagestyle{fancy}
\fancyhf{}
\lhead{Introduction to Deep Learning and Practice}
\rhead{Deep Learning: Answers}
\cfoot{\thepage}
\renewcommand{\headrulewidth}{0.4pt}
\setlength{\headheight}{14pt}

\hypersetup{
  colorlinks=true,
  linkcolor=deepblue,
  urlcolor=deepblue
}

% =================================================

% source tag of a problem
\newcommand{\src}[1]{{\unskip\nobreak\hfil\penalty50\hskip1em\hbox{}\nobreak\hfill\mbox{\small\color{black!50}[#1]}}}
% a problem block
\newlist{problems}{enumerate}{1}
\setlist[problems]{label=\textbf{\arabic*.},leftmargin=2.2em,itemsep=0.2em}
% =================================================
% a solution block
\newtcolorbox{sol}{breakable,enhanced,colback=softgray,colframe=black!35,boxrule=0.4pt,arc=1mm,
  left=2mm,right=2mm,top=1mm,bottom=1mm,before upper={\textbf{풀이.}\ }}
\setlist[problems]{label=\textbf{\arabic*.},leftmargin=2.2em,itemsep=0.9em}
\begin{document}

\begin{center}
{\LARGE\bfseries 연습문제 정답과 풀이}\\[0.35em]
{\large 딥러닝 $\cdot$ 딥러닝개론및연습 2026년 2학기}
\end{center}

강의노트 끝의 연습문제 10개의 정답과 풀이이다. 문제를 다시 적고 그 아래에 풀이를 쓴다.
7번 (c)와 10번의 숫자는 이 컴퓨터에서 코드를 실제로 실행한 결과이다.

\begin{problems}

% ---------------------------------------------------------------- 1
\item DeepConvNet의 Conv4의 패딩을 $2$가 아니라 $1$로 바꾸면 Conv4 다음 풀링의 출력, Conv6 다음 풀링의 출력,
Flatten의 출력 크기는 각각 얼마가 되는가? 첫 번째 Affine 계층의 매개변수 수는 어떻게 바뀌는가?

\begin{sol}
패딩 $1$인 $3\times3$ 합성곱은 크기를 바꾸지 않으므로 Conv4의 출력은 $(32,14,14)$이다.
\begin{center}
\begin{tabular}{lcc}
\toprule
 & 원래 (Conv4 패딩 $2$) & Conv4 패딩 $1$ \\
\midrule
Conv4 & $(32,16,16)$ & $(32,14,14)$ \\
Conv4 다음 풀링 & $(32,8,8)$ & $\mathbf{(32,7,7)}$ \\
Conv5, Conv6 & $(64,8,8)$ & $(64,7,7)$ \\
Conv6 다음 풀링 & $(64,4,4)$ & $\mathbf{(64,3,3)}$ ($\lfloor7/2\rfloor=3$) \\
Flatten & $1024$ & $\mathbf{576}$ \\
첫 번째 Affine & $1024\cdot50+50=51250$ & $576\cdot50+50=\mathbf{28850}$ \\
\bottomrule
\end{tabular}
\end{center}
첫 번째 Affine의 매개변수가 $22400$개 줄어 전체는 $123552\to101152$개가 된다.
대신 마지막 풀링에서 $7\times7$의 마지막 행과 열 (가장자리 한 줄)이 버려진다.
Conv4의 패딩을 $2$로 한 것은 $16\to8\to4$로 나누어떨어지게 하여 이 손실을 피하려는 것이다.
\end{sol}

% ---------------------------------------------------------------- 2
\item 입력과 출력의 채널이 모두 $C$개일 때, $7\times7$ 합성곱 계층 하나와 $3\times3$ 합성곱 계층 세 개 (패딩 $1$, 스트라이드 $1$)를 비교하라.
(a) 출력 원소 하나의 수용 영역이 같음을 보여라. (b) 편향을 빼고 매개변수 수를 비교하라. (c) $C=64$일 때 각각의 값을 구하라.

\begin{sol}
(a) $3\times3$ 합성곱의 출력 하나는 바로 앞 층의 $3\times3$ 영역을 본다.
세 번째 층의 출력 하나 $\leftarrow$ 둘째 층의 $3\times3$ $\leftarrow$ 첫째 층의 $5\times5$ $\leftarrow$ 입력의 $7\times7$이다.
한 층을 지날 때마다 한 변이 $2$씩 늘어 $3\to5\to7$이므로, 입력에서 보면 수용 영역은 $7\times7$로 $7\times7$ 합성곱 하나와 같다.

(b) $7\times7$ 하나: $7\cdot7\cdot C\cdot C=49C^2$. \quad $3\times3$ 세 개: $3\cdot(3\cdot3\cdot C\cdot C)=27C^2$.
$3\times3$ 세 개가 약 $27/49\approx55\%$의 매개변수로 같은 영역을 본다. 게다가 중간에 ReLU가 두 번 더 들어가 비선형성이 늘어난다.

(c) $C=64$: $49\cdot4096=\mathbf{200704}$개 vs $27\cdot4096=\mathbf{110592}$개.
\end{sol}

% ---------------------------------------------------------------- 3
\item VGG16의 합성곱 계층의 매개변수 수를 구하라. (첫 블록: $3\to64$, $64\to64$; 둘째: $64\to128$, $128\to128$;
셋째: $128\to256$, $256\to256$ 두 번; 넷째: $256\to512$, $512\to512$ 두 번; 다섯째: $512\to512$ 세 번. 모두 $3\times3$, 편향 포함.)
완전연결 계층 세 개 ($7\cdot7\cdot512\to4096\to4096\to1000$)의 매개변수 수와 비교하라.

\begin{sol}
$c_{\text{in}}\to c_{\text{out}}$인 $3\times3$ 합성곱의 매개변수는 $9\,c_{\text{in}}c_{\text{out}}+c_{\text{out}}$이다.
\begin{center}
\begin{tabular}{lr}
\toprule
블록 & 매개변수 수 \\
\midrule
1: $3\to64$, $64\to64$ & $1792+36928=38720$ \\
2: $64\to128$, $128\to128$ & $73856+147584=221440$ \\
3: $128\to256$, $256\to256$ ($\times2$) & $295168+590080\cdot2=1475328$ \\
4: $256\to512$, $512\to512$ ($\times2$) & $1180160+2359808\cdot2=5899776$ \\
5: $512\to512$ ($\times3$) & $2359808\cdot3=7079424$ \\
\midrule
합성곱 합계 & $\mathbf{14714688}$ (약 $1470$만) \\
\bottomrule
\end{tabular}
\end{center}
완전연결 계층: $25088\cdot4096+4096=102764544$, $4096\cdot4096+4096=16781312$, $4096\cdot1000+1000=4097000$, 합계 $\mathbf{123642856}$ (약 $1$억 $2400$만).

전체 $138357544$개 (약 $1.38$억) 중 \textbf{완전연결 계층이 약 $89\%$}이고, 그중 첫 번째 ($7\cdot7\cdot512\to4096$) 하나가 약 $1$억 개이다.
합성곱 계층은 층이 $13$개나 되지만 매개변수는 $11\%$뿐이다. 반면 계산량 (곱셈 횟수)은 큰 특징 맵 위를 움직이는 합성곱 계층이 대부분을 차지한다.
그래서 나중의 network (GoogLeNet, ResNet)는 큰 완전연결 계층 대신 전역 평균 풀링을 써서 매개변수를 크게 줄였다.
\end{sol}

% ---------------------------------------------------------------- 4
\item 입력이 $192$채널이고 $3\times3$ 필터 $128$개를 적용한다. (a) 바로 적용할 때와 (b) $1\times1$ 합성곱으로 $96$채널로 줄인 뒤 적용할 때,
편향을 빼고 매개변수 수를 비교하라. 입력의 공간 크기가 $28\times28$ (패딩 $1$)이면 곱셈 횟수는 각각 얼마인가?

\begin{sol}
(a) $3\cdot3\cdot192\cdot128=\mathbf{221184}$개.\\
(b) $1\times1$: $192\cdot96=18432$, $3\times3$: $3\cdot3\cdot96\cdot128=110592$, 합계 $\mathbf{129024}$개 (약 $58\%$).

패딩 $1$이므로 출력도 $28\times28=784$개 위치이고, 위치마다 필터의 가중치 수만큼 곱한다.
\[
\text{(a) } 221184\times784=\mathbf{173408256}\ (\text{약 }1.73\text{억}),\qquad
\text{(b) } 129024\times784=\mathbf{101154816}\ (\text{약 }1.01\text{억}).
\]
$1\times1$ 합성곱은 공간은 그대로 두고 채널만 섞어 줄이는 ``병목'' 역할을 한다. 그래서 같은 출력을 약 $42\%$ 적은 계산으로 만든다.
GoogLeNet의 인셉션 구조와 ResNet의 병목 블록이 이 방법을 쓴다.
\end{sol}

% ---------------------------------------------------------------- 5
\item 스칼라 블록 $30$개를 생각하자. 일반 network는 $h_{l}=w_l h_{l-1}$, ResNet은 $h_{l}=h_{l-1}+w_l h_{l-1}$이다.
모든 $w_l=0.5$일 때 $\pd{h_{30}}{h_0}$를 각각 구하라. $w_l$이 모두 $0$이면 어떻게 되는가?
이로부터 ``필요 없는 블록은 $F(x)=0$을 배우면 된다''는 말을 설명하라.

\begin{sol}
연쇄법칙으로 $\pd{h_{30}}{h_0}=\prod_{l=1}^{30}\pd{h_l}{h_{l-1}}$이다.
\begin{itemize}[leftmargin=1.5em,itemsep=2pt]
\item 일반: $\pd{h_l}{h_{l-1}}=w_l=0.5$이므로 $0.5^{30}\approx\mathbf{9.3\times10^{-10}}$. 기울기가 사라진다.
\item ResNet: $\pd{h_l}{h_{l-1}}=1+w_l=1.5$이므로 $1.5^{30}\approx\mathbf{1.9\times10^{5}}$.
\end{itemize}
(ResNet에서 기울기가 커진 것은 $w_l=0.5$로 모든 블록이 신호를 키우기 때문이다.
실제 ResNet은 배치 정규화와 작은 초깃값으로 $F$를 작게 두어 이런 폭발을 막는다. 핵심은 곱해지는 수가 $w_l$이 아니라 $1+w_l$이라는 점이다.)

$w_l$이 모두 $0$이면:
\begin{itemize}[leftmargin=1.5em,itemsep=2pt]
\item 일반: $h_{30}=0$이고 $\pd{h_{30}}{h_0}=0$이다. 입력의 정보도 기울기도 전혀 전달되지 않는다.
\item ResNet: $h_l=h_{l-1}$이므로 $h_{30}=h_0$이고 $\pd{h_{30}}{h_0}=1$이다. 입력과 기울기가 그대로 전달된다.
\end{itemize}

\textbf{해석.} ResNet의 블록은 $y=x+F(x)$이다. $F(x)=0$이면 블록이 항등함수가 되어 아무 일도 하지 않는다.
그래서 층을 더 쌓아도 쓸모없는 블록은 $F$의 가중치를 $0$ 근처로 배우기만 하면 원래의 얕은 network와 같아진다.
일반 network는 ``아무것도 하지 않는'' 층 $h_l=h_{l-1}$을 만들려면 가중치가 정확히 항등함수 ($w=1$)가 되어야 하는데, 이는 $0$ 근처에서 시작하는 학습에게 어렵다.
즉 skip connection이 있으면 깊게 쌓아도 적어도 얕은 network보다 나빠지기 어렵고, 기울기는 덧셈 경로를 따라 그대로 앞까지 전달된다.
\end{sol}

% ---------------------------------------------------------------- 6
\item 미니배치 $B=6$개의 손실이 $\ell_1,\ldots,\ell_6$이고 $L=\frac16\sum\ell_n$이다.
GPU 두 개가 각각 $L_1=\frac13(\ell_1+\ell_2+\ell_3)$, $L_2=\frac13(\ell_4+\ell_5+\ell_6)$의 기울기를 구했다.
$\grad L=\frac12(\grad L_1+\grad L_2)$임을 보여라. 조각의 크기가 $4$개와 $2$개로 다르면 어떻게 합쳐야 하는가?

\begin{sol}
미분은 선형이므로
\[
\grad L=\frac16\sum_{n=1}^{6}\grad\ell_n
=\frac12\Bigl(\frac13\sum_{n=1}^{3}\grad\ell_n+\frac13\sum_{n=4}^{6}\grad\ell_n\Bigr)
=\frac12\bigl(\grad L_1+\grad L_2\bigr).
\]
조각이 $4$개와 $2$개이면 $L_1=\frac14(\ell_1+\cdots+\ell_4)$, $L_2=\frac12(\ell_5+\ell_6)$이고
\[
\grad L=\frac16\Bigl(\sum_{n=1}^{4}\grad\ell_n+\sum_{n=5}^{6}\grad\ell_n\Bigr)
=\frac16\bigl(4\grad L_1+2\grad L_2\bigr)=\frac23\grad L_1+\frac13\grad L_2 .
\]
즉 조각의 크기에 비례하는 \textbf{가중 평균}으로 합쳐야 한다. 단순 평균 $\frac12(\grad L_1+\grad L_2)$를 쓰면 작은 조각의 데이터 $2$개가 지나치게 큰 비중을 갖는다.
(지난주 코드에서 \texttt{loss.item() * len(yb)}를 더한 뒤 전체 개수로 나눈 것과 같은 이유이다.)
\end{sol}

% ---------------------------------------------------------------- 7
\item \texttt{float16}은 가수가 $10$비트이다. (a) $1$ 바로 다음으로 나타낼 수 있는 수는? (b) $2048$ 바로 다음 수는?
(c) numpy로 \texttt{np.float16(2048) + np.float16(1)}을 계산하고 결과를 설명하라.
(d) 학습 중 아주 작은 기울기 (예: $10^{-8}$)를 \texttt{float16}으로 바꾸면 어떻게 되는가?

\begin{sol}
float16의 수는 $\pm(1.f)_2\times2^{e}$로, 가수 $f$가 $10$비트이다. 그래서 $[2^e,\,2^{e+1})$ 구간에서 나타낼 수 있는 수의 간격은 $2^{e-10}$이다.
\begin{itemize}[leftmargin=1.5em,itemsep=2pt]
\item (a) $1=2^0$ 근처의 간격은 $2^{-10}$이므로 다음 수는 $1+2^{-10}=\mathbf{1.0009765625}$이다.
(\texttt{np.nextafter(np.float16(1), np.float16(2))}의 결과와 같다.)
\item (b) $2048=2^{11}$ 근처의 간격은 $2^{11-10}=2$이므로 다음 수는 $\mathbf{2050}$이다 (numpy로도 $2050$).
\item (c) 실행 결과는 \texttt{2048.0} (\texttt{float16})이다.
참값 $2049$는 나타낼 수 없고, $2048$과 $2050$의 정확히 가운데이다. 이때는 가수의 마지막 비트가 짝수인 쪽으로 반올림하므로 $2048$이 된다.
즉 $1$을 더했는데 아무 변화가 없다. float16으로 큰 값에 작은 값을 계속 더하면 (예: 기울기를 누적하면) 더한 것이 사라질 수 있다.
\item (d) float16이 나타낼 수 있는 가장 작은 양수는 $2^{-24}\approx5.96\times10^{-8}$이다 (정규화된 수의 최솟값은 $2^{-14}\approx6.1\times10^{-5}$이고, 그보다 작은 수는 정밀도를 잃으며 표현된다).
$10^{-8}$은 그 절반보다 작으므로 \texttt{np.float16(1e-8)}은 $\mathbf{0}$이 된다 (underflow, 실제로 확인).
작은 기울기가 $0$이 되면 그 매개변수는 갱신되지 않는다.
그래서 혼합 정밀도 학습에서는 손실에 큰 수 (예: $2^{16}$)를 곱해 기울기를 키운 뒤 float16으로 계산하고, 갱신 직전에 다시 나누는 \textbf{loss scaling}을 쓴다.
매개변수 자체의 갱신은 float32로 하여 (c)의 문제를 피한다.
\end{itemize}
\end{sol}

% ---------------------------------------------------------------- 8
\item 정답 사각형이 $(x,y)=(0,0)$부터 $(4,4)$까지이고, 예측한 사각형이 $(1,1)$부터 $(5,5)$까지이다. IoU를 구하라.
IoU $\ge0.5$를 맞은 것으로 볼 때, 예측 사각형을 정답과 같은 크기로 오른쪽으로만 $d$만큼 옮긴다면 $d$가 얼마까지 맞은 것인가?

\begin{sol}
두 사각형의 넓이는 모두 $16$이다. 겹치는 부분은 $(1,1)$부터 $(4,4)$까지로 넓이 $3\cdot3=9$이다.
\[
\text{IoU}=\frac{\text{교집합}}{\text{합집합}}=\frac{9}{16+16-9}=\frac{9}{23}\approx\mathbf{0.391}.
\]
$0.5$보다 작으므로 맞은 것으로 보지 않는다.

예측 사각형이 $(d,0)$부터 $(4+d,4)$까지이면 ($0\le d\le4$) 교집합은 $4(4-d)$, 합집합은 $32-4(4-d)=16+4d$이다.
\[
\text{IoU}=\frac{16-4d}{16+4d}\ge\frac12
\iff 32-8d\ge16+4d
\iff d\le\frac43\approx1.33 .
\]
즉 한 변이 $4$인 사각형을 $\mathbf{4/3}$ (한 변의 $\frac13$)만큼까지 옮겨도 맞은 것이다.
처음의 예측은 가로와 세로로 $1$씩 어긋나서, 한 방향으로만 $1$ 옮긴 경우 ($\text{IoU}=\frac{12}{20}=0.6$)보다 IoU가 작다.
\end{sol}

% ---------------------------------------------------------------- 9
\item MNIST의 데이터 확장으로 다음 중 무엇이 적절한가? 이유와 함께 답하라.
(a) $\pm10^\circ$ 회전 (b) 좌우 뒤집기 (c) $180^\circ$ 회전 (d) $\pm2$픽셀 이동 (e) 약간의 굵기 변화
(f) 고양이·개 사진 분류라면 (b)는 어떤가?

\begin{sol}
기준은 ``변형해도 \textbf{정답이 바뀌지 않고}, 실제 데이터에 나올 법한가''이다.
\begin{itemize}[leftmargin=1.5em,itemsep=2pt]
\item (a) \textbf{적절}. 사람마다 글씨가 조금씩 기울어진다. 노트의 \texttt{augment}도 $\pm10^\circ$를 쓴다.
\item (b) \textbf{부적절}. 숫자를 좌우로 뒤집으면 실제로 쓰지 않는 모양이 된다 ($2$, $3$, $7$ 등). 뒤집은 $2$는 $5$와 비슷해 보이기도 한다.
\item (c) \textbf{부적절}. $6$이 $9$가 되고 $9$가 $6$이 되어 정답이 바뀐다.
\item (d) \textbf{적절}. 숫자가 칸의 가운데에서 조금 비켜 있어도 같은 숫자이다. 합성곱과 풀링이 어느 정도 이동에 강하지만 데이터로도 알려 준다.
\item (e) \textbf{적절}. 펜에 따라 획의 굵기가 다르다. (너무 굵게 하면 $8$의 구멍이 메워지는 등 모양이 바뀌므로 약간만 한다.)
\item (f) \textbf{적절}. 좌우를 뒤집은 고양이도 고양이이고, 사진은 어느 쪽을 보고 찍혀도 된다. 사진 분류에서 가장 흔히 쓰는 확장이다.
(위아래 뒤집기는 거꾸로 선 고양이 사진이 드물어서 보통 쓰지 않는다.)
\end{itemize}
\end{sol}

% ---------------------------------------------------------------- 10
\item \texttt{mnist\_deep\_torch.py --n\_train 5000}과 \texttt{--n\_train 5000 --augment}로 학습하고 test 정확도를 비교하라.
1.2절의 표와 함께, 훈련 데이터의 양에 따라 데이터 확장의 효과가 어떻게 달라지는지 설명하라.

\begin{sol}
노트의 코드를 그대로 실행하였다 (기본값 $20$ epoch, Adam, 미니배치 $100$, 같은 초깃값 씨앗). 노트의 표에 $5000$장의 결과를 더하면 다음과 같다.
\begin{center}
\begin{tabular}{@{}lcccc@{}}
\toprule
훈련 데이터 (epoch) & 확장 없음 & 확장 & 차이 & 틀린 test 이미지 \\
\midrule
$1000$장 ($60$, 노트) & $95.38\%$ & $96.76\%$ & $+1.38\%$p & $462\to324$ ($-30\%$) \\
\textbf{$5000$장 ($20$)} & $\mathbf{97.74\%}$ & $\mathbf{98.59\%}$ & $+0.85\%$p & $226\to141$ ($-38\%$) \\
$55000$장 ($20$, 노트) & $99.49\%$ & $99.50\%$ & $+0.01\%$p & $51\to50$ ($-2\%$) \\
\bottomrule
\end{tabular}
\end{center}
($5000$장일 때 마지막 epoch의 검증 정확도는 확장 없음 $97.98\%$, 확장 $98.34\%$였다.)

\textbf{해석.}
\begin{itemize}[leftmargin=1.5em,itemsep=2pt]
\item 훈련 데이터가 적을수록 정확도의 상승폭이 크다 ($+1.38\to+0.85\to+0.01\%$p).
데이터가 적으면 network가 훈련 이미지를 외우기 쉽고 (과대적합), 회전·이동한 새 이미지가 ``보지 못한 변형''을 채워 준다.
$55000$장이면 여러 기울기와 위치의 숫자를 이미 충분히 보았으므로 확장이 새로 알려 주는 것이 거의 없다.
\item $5000$장에서는 틀린 이미지가 $38\%$ 줄어, 비율로 보면 $1000$장 ($30\%$)과 비슷하거나 더 크다.
\item $5000$장 $+$ 확장 ($98.59\%$)은 확장 없이 $55000$장을 쓴 결과 ($99.49\%$)에는 못 미친다. 확장은 진짜 데이터를 대신하지 못하지만, 적은 데이터에서 그 차이를 상당히 줄여 준다.
\item 주의: 각 조건을 씨앗 하나로 한 번씩만 실행하였다. $1000$장은 $60$ epoch, $5000$장은 $20$ epoch로 학습한 양도 다르다.
$55000$장에서의 $0.01\%$p (이미지 $1$장)처럼 작은 차이는 씨앗만 바꾸어도 생길 수 있다. 그러나 $5000$장에서의 $85$장 차이는 그보다 훨씬 크다.
\end{itemize}
\end{sol}

\end{problems}

\end{document}
